\begin{center}
	\textbf{Abstract}
\end{center}


The liquid, outer core of the Earth undergoes vigorous and turbulent convection. 
The underlying dynamics are poorly understood and, in the absence of sufficient observational data, numerical and mathematical models are used to investigate processes in the core. 
In this dissertation we present three numerical investigations of fluid mechanical processes relevant to the outer core of the Earth.
We discuss the scales and temporal behavior of the turbulence which underpins these processes, and the ability of this turbulence to transport heat, momentum, and angular momentum from the interior of the planet out to the solid mantle. 
In our first study we investigate the small scales of convective turbulence, and demonstrate how previous work may have underestimated the size of the smallest eddies relevant to planetary interiors. We propose a new theory which suggests that molecular viscosity plays a significant role in the dynamics of planetary scale flows.
In our next two studies we zoom out, and discuss how large scale, planetary features interact with a convecting core. In particular, we investigate the role of topography at the core mantle boundary. Whether topographic interactions at the core mantle boundary can account for known variations in the length of day is an outstanding question, and here we present the first direct numerical simulations of a turbulent, rapidly rotating spherical shell with boundary topography. We investigate local and global scale topography and determine that the pressure torques exerted by the core on the core mantle boundary are sufficiently large to account for variations in Earth's length of day.
